% !PRINTED 195
\weberpage{195}
\begin{center}
\S.\,51.

\gesperrt{Endliche Gruppen linearer gebrochener Substitutionen.}

\gesperrt{Pole der Gruppen.}
\end{center}

Aus den bisherigen Entwickelungen ergiebt sich, dass, wenn
alle endlichen Gruppen linearer gebrochener Substitutionen gefunden
 sind, daraus ohne Weiteres alle endlichen Gruppen
eigentlich orthogonaler ternärer Substitutionen und alle endlichen
 Gruppen binärer linearer Substitutionen mit der Determinante
 1 gefunden werden können.

Ausserdem giebt es noch endliche Gruppen, die neben den
eigentlichen auch uneigentlich orthogonale Substitutionen enthalten
.

Diese Gruppen, auf die wir später zurückkommen, können
nicht aus den Gruppen linearer gebrochener Substitutionen allein
abgeleitet werden, weil sich bei den gebrochenen Substitutionen
der Unterschied zwischen eigentlich und uneigentlich orthogonalen
 Substitutionen verwischt.

Wir suchen also jetzt zunächst \gesperrt{alle endlichen Gruppen
linearer gebrochener Substitutionen} zu ermitteln $^{1)}$.

Wir bezeichnen mit $G$ eine solche Gruppe vom Grade $n$,
und mit
\weq{1}{x,\ \Theta_1(x),\ \Theta_2(x),\ \ldots,\ \Theta_{n-1}(x)}
ihre Elemente, worin die $\Theta$ Symbole für lineare Functionen
\weq{2}{\Theta(x) = \frac{\alpha x + \beta}{\gamma x + \delta}}
sind, die auch mit
\weq{3}{\Theta = \begin{pmatrix}\alpha, & \beta\\ \gamma, & \delta\end{pmatrix}}
bezeichnet werden.

Aus jeder Gruppe von der Form (1) können wir nach \S.\,37
eine ganze Schaar isomorpher Gruppen ableiten, wenn wir mit
$L$ eine willkürliche lineare Substitution bezeichnen:
\weq{4}{1,\ L\Theta_1L^{-1},\ L\Theta_2L^{-1},\ \ldots,\ L\Theta_{n-1}L^{-1},}
%CONT>>

\par\noindent\rule{4cm}{0.4pt}\par
$^{1)}$ Ueber die Theorie dieser Gruppen ist besonders zu vergleichen:
\gesperrt{Gordan}, Ueber endliche Gruppen linearer Transformationen einer
Veränderlichen. Mathematische Annalen, Bd. XII, S. 23 (1877);
\gesperrt{Klein}, Vorlesungen über das Ikosaëder (Leipzig 1884).
